\documentclass[10pt]{article}
\usepackage[margin=20mm]{geometry}
\usepackage{amsmath,amssymb,mathtools,fontspec,unicode-math,xurl,hyperref}
\IfFontExistsTF{DejaVu Serif}{\setmainfont{DejaVu Serif}}{\setmainfont{Latin Modern Roman}}
\setmathfont{latinmodern-math.otf}
\setlength{\parindent}{0pt}\setlength{\parskip}{6pt}\sloppy
\title{Scientific initial review M25-0009
De Rham algebras of closed quasiregularly elliptic manifolds are Euclidean}
\author{}
\date{26 October 2026}
\begin{document}
\maketitle
=3em Object and disposition. This report reviews Susanna Heikkilä and Pekka Pankka’s arXiv:2302.11440v2, dated 7 December 2023, 25 PDF pages. It records two Major issues in the written argument or its advertised scope and one Minor issue in an example. One Major issue has a complete local repair supplied below. The other concerns a missing implication for an arbitrary given smooth structure, not a constructed counterexample to the classification. Neither finding disproves the principal algebra-embedding theorem.\par
The corresponding journal identity is Annals of Mathematics 201 (2025), 459–488, DOI 10.4007/annals.2025.201.2.3. The journal Version of Record was not obtained or verified as textually identical. All findings in this report are confined to the supplied v2.\par
\section*{Scope, versions, and method}
All 25 PDF pages were read, including definitions, all proof sections, examples, appendices, and references. Source PDF pages 2, 4, 14, 17, 19, 22, 23, and 24 were additionally inspected as images. The supplied 1,437-line author TeX file was inspected statically; it was not executed or compiled. The manifest’s seven file bindings were actually hashed and matched. Independent calculations below use only the reviewing assistant’s own Python standard-library script.\par
The review followed the chain from the rescaling/hunting argument to bounded normalized pullbacks, compact homotopy primitives, weak wedge convergence, the top-degree limit measure, multiplicativity, Poincaré-duality injectivity, and pointwise evaluation. It also checked the bounded-energy alternative, the singular sphere example, intersection-form bounds, and the PL-to-quasiregular interface.\par
The arXiv history inspected during this review lists v1 on 22 February 2023 and v2 on 7 December 2023. The publisher metadata reports a later revision on 20 November 2024 and publication online on 12 March 2025; its displayed acceptance date, 16 May 2024, precedes that revision and is retained as displayed rather than silently reordered. The official Helsinki research portal lists an accepted manuscript named Heikkila-Pankka\_QR-cohomology-revised-Nov-2024. A bounded retrieval received HTTP 403. A publicly indexed Helda reprint was also found, but opening it encountered an access-denied challenge. That barrier was not bypassed. An indexed excerpt cannot establish the content of the entire later manuscript. No claim is made that any issue below persists in, or was repaired by, either later text. Bounded publisher/arXiv searches did not identify an erratum; absence of a search hit is not proof of absence.\par
\section*{F01 — Major: division by a possibly zero top-degree pairing}
Equation reference: eq:bad. Exact location. Proposition 5.1, PDF page 14, the displayed assertion immediately after the case \(k_1+k_2=n\) begins; author TeX lines 860–870, especially line 862. For arbitrary classes \(c_1,c_2\) in complementary positive degrees, the proof selects \(\tau\) such that \[
 \operatorname{vol}(N)h(c_1)\wedge h(c_2)
 = I\bigl(\operatorname{vol}_N+d\tau\bigr),\qquad
 I=\int_N c_1\wedge c_2 .\] Here \(h(c)\) is the harmonic representative. No nonzero-pairing hypothesis is imposed in this proposition. If \(I=0\), this would require the harmonic representatives to have identically zero wedge, which need not hold.\par
\subsection*{A counterexample satisfying the actual surrounding setup}
Let \(N=\mathbb R^3/\mathbb Z^3\) with coordinates \(x,y,z\) modulo one, oriented by \(dx\wedge dy\wedge dz\). Put \(a(z)=\frac12\sin(2\pi z)\) and choose the smooth metric whose inverse matrix is \[(g^{ij})=
 \begin{pmatrix}
 1&a&0\\ a&1&0\\0&0&(1-a^2)^{-1}
 \end{pmatrix},
 \qquad
 (g_{ij})=
 \begin{pmatrix}
 (1-a^2)^{-1}&-a(1-a^2)^{-1}&0\\
 -a(1-a^2)^{-1}&(1-a^2)^{-1}&0\\
 0&0&1-a^2
 \end{pmatrix}.\] Both matrices are positive definite, \(\det g=1\), and the largest eigenvalue of \(g\) is at most two because \(|a|\le1/2\). Hence \(\operatorname{vol}(N)=1\) and \(\operatorname{vol}_N=dx\wedge dy\wedge dz\). Directly from the defining identity for the Hodge star, \[*dx=dy\wedge dz+a\,dz\wedge dx,\qquad
 *dy=a\,dy\wedge dz+dz\wedge dx .\] These two forms are closed: each differentiated coefficient already multiplies a factor \(dz\). Thus \(dx\) and \(dy\) are closed and coclosed. Set \[\alpha=dx,\qquad \beta=*dy.\] The form \(\beta\) is closed, and \(*\beta=dy\) is closed, so \(\beta\) is also harmonic. With \(c_1=[\alpha]\) and \(c_2=[\beta]\), uniqueness of harmonic representatives gives \(h(c_1)=\alpha\), \(h(c_2)=\beta\). However, \[\alpha\wedge\beta=a(z)\operatorname{vol}_N,\qquad
 \int_N\alpha\wedge\beta=0,\qquad
 a(1/4)=1/2.\] The left side of [eq:bad] (\#eq:bad) is nonzero and its right side is zero for every \(\tau\). For clarity, the nonzero wedge is exact: \[\alpha\wedge\beta=d\eta_0,\qquad
 \eta_0=-\frac{\cos(2\pi z)}{4\pi}\,dx\wedge dy .\] Exactness is weaker than pointwise vanishing.\par
This example is not placed on a target excluded by the unbounded-area hypothesis. For every positive integer \(j\) take \[f_j:B_2^3\longrightarrow N,\qquad u\longmapsto ju\pmod{\mathbb Z^3}.\] Its Jacobian is \(j^3\), and \(\|Df_j\|^3\le2\sqrt2\,j^3\). Writing \(\omega_3=\operatorname{vol}_{\mathbb R^3}(B^3)\), one obtains \[A(f_j)=\omega_3j^3\longrightarrow\infty,\qquad
 \int_{B_2^3}f_j^*\operatorname{vol}_N=8A(f_j).\] Thus \(f_j\in\mathcal F_{2\sqrt2,8}(N)\) with fixed constants. The target is closed, connected, oriented, and not a rational homology sphere. The normalized pullbacks of a finite harmonic basis are bounded in the required \(L^{3/k}\) spaces. Weak compactness and compactness of the Poincaré operator give a Sobolev–Poincaré subsequence, exactly as in Lemma 4.1, independently of the invalid assertion [eq:bad] (\#eq:bad). In the two specified degrees the normalized forms are explicitly \[f_j^\#(c_1)=\omega_3^{-1/3}dx,\qquad
 f_j^\#(c_2)=\omega_3^{-2/3}
       \bigl(a(jz)\,dy\wedge dz+dz\wedge dx\bigr).\] Periodic averaging makes the oscillatory coefficient converge weakly to zero. The limiting product is zero, as the proposition’s desired conclusion requires.\par
\subsection*{Repair, counterevidence, and core impact}
Equation reference: eq:repair. The correct top-cohomology decomposition, valid also when \(I=0\), is \[
 h(c_1)\wedge h(c_2)=\frac{I}{\operatorname{vol}(N)}\operatorname{vol}_N+d\eta .\] Its difference has integral zero and hence is exact on the closed connected oriented \(N\). Applying normalized pullback, testing against a compactly supported smooth function, using Proposition 4.2 for weak wedge convergence, Lemma 3.2 for the vanishing exact term, and Theorem 4.4 for the volume-form limit gives \[d\widehat L(c_1)\wedge d\widehat L(c_2)
  =\frac{I}{\operatorname{vol}(N)}\,\nu_{d\widehat L}
  =\Phi_{d\widehat L}(c_1\wedge c_2).\] This supplies the missing zero-pairing case without division by \(I\). The earlier use of a similar formula in Theorem 4.4, page 12, explicitly chooses \(I\ne0\) by Poincaré duality and is not defective for this reason.\par
The invalid step lies on Proposition 5.1 \(\to\) Theorem 1.7 \(\to\) Theorem 1.3 \(\to\) Theorem 1.1. Under standard 0.7.1 it is Major because a core written derivation fails under its own hypotheses, although the local repair is complete and no main-theorem counterexample results. The repaired statement is adopted locally in this report; that adoption does not retroactively change the reviewed source.\par
\section*{F02 — Major: the sufficiency argument changes the smooth/PL scope}
Exact location. Corollary 1.2 and the two preceding paragraphs on PDF page 2, author TeX lines 398–415; compare Proposition B.1 on page 23, lines 1206–1213. The definition on page 1 is for a given Riemannian manifold. The corollary asserts an if-and-only-if condition on that manifold using only its homeomorphism type. Its sufficiency argument, however, constructs torus branched covers for the displayed connected sums with their standard PL structures compatible with the usual smooth structures.\par
Dependency verification. The inspected primary version of Piergallini–Zuddas [PZ], Theorem 1.3, constructs branched covers onto the standard connected-sum models. Remark 1.4 explains the torus consequence and expressly leaves the exotic counterparts outside that conclusion. Corollary 1.5 constructs covers from an exotic counterpart to a standard model, not covers in the reverse direction. This describes the scope of that primary source; it is not a claim about the current status of every exotic case.\par
The missing implication. Let \(X\) be a manifold with its given smooth structure and let \(N_0\) be a standard model on the list. The supplied premises give \[X\cong_{\mathrm{TOP}}N_0,\qquad
 \mathbb R^4\longrightarrow\mathbb T^4\longrightarrow N_0
 \quad\hbox{quasiregular}.\] To conclude that \(X\) itself is quasiregularly elliptic by this construction requires a suitable quasiregular transfer to \(X\), or a separate branched-cover construction for its compatible PL structure. An arbitrary homeomorphism \(N_0\to X\) has no stated Sobolev or bounded-distortion control. Proposition B.1 does not supply it: that proposition requires a PL map and PL structures compatible with both given smooth structures. A topological classification alone does not verify those hypotheses. Nor can a branched covering \(X\to N_0\) simply be inverted and composed.\par
Consequently, the proof establishes a precise weaker classification: every closed simply connected quasiregularly elliptic smooth four-manifold has one of the listed homeomorphism types, and every listed type has a quasiregularly elliptic standard smooth representative. It does not establish, by the cited steps, the stronger sufficiency for every given smooth structure on such a type. No existence assertion about a particular exotic structure in this small range, and no nonelliptic exotic example, is needed or claimed here.\par
Counterevidence and resolution. The standard models really do have the required torus covers. The necessary intersection-form bounds \(\beta_2^\pm\le3\) are independently checked below. Changing the smooth Riemannian metric on a fixed closed smooth manifold is harmless: uniform equivalence of the metrics controls distortion. These facts do not address a change of compatible PL/smooth structure. The repair is either to state the weaker homeomorphism-type classification explicitly or to provide the missing theorem for the stronger quantifier. The abstract’s advertised classification and Corollary 1.2 are a main application, so this demonstrated scope gap is Major. The truth of the stronger formulation remains unresolved by this review; it is not labeled false.\par
\section*{F03 — Minor: the zoom in Example 7.3 includes a transition annulus}
Exact location. Example 7.3, PDF page 17, author TeX lines 992–996. First, the prescribed balls \(B(2^j e_1,j)\) are not mutually disjoint for all \(j\ge1\): the pairs \((1,2)\) and \((2,3)\) overlap. Starting at \(j=3\) repairs this finite-prefix defect.\par
Even after that repair, the positive zoom \(f_{a_j,r_j/2}\) is defined on \(B_2^n\). It sees the whole ball of radius \(r_j\), not only the inner ball on which \(h_Q\) is the fixed rotation \(Q\). In normalized coordinates its map is \[H(x)=
 \begin{cases}
 Qx,& |x|\le1,\\
 Q_{2-|x|}x,&1\le |x|<2.
 \end{cases}\] This is independent of \(j\). Taking \(Q_t=\exp(tA)\) with \(A\) skew-symmetric, the derivative of a radial twist \(x\mapsto Q_{t(|x|)}x\) has determinant \[\det Q_t\left(1+\frac{t'(|x|)}{|x|}x\cdot Ax\right)=1.\] Consequently normalization does not remove the annular derivative: its degree-one pullback is \(\operatorname{vol}(B^n)^{-1/n}dH\), not the constant form specified by \(Q^*L_\pi\) on all of \(B_2^n\).\par
For an explicit check, take \(n=2\), \(Q\) a rotation by \(\theta=\pi/2\), and \(x=(3/2,0)\). Then \[DH(x)=R_{\pi/4}
 \begin{pmatrix}1&0\\-3\pi/4&1\end{pmatrix},\] which is different from \(Q=R_{\pi/2}\). Continuity on the open annulus gives a positive-measure discrepancy, not just a single exceptional point. The independent numerical check gives a maximum entrywise difference approximately \(2.37319\).\par
Use the positive zoom \(f_{a_j,r_j/4}\) instead. Its full domain \(B_2^n\) then lies in the inner rotation ball. For a particularly transparent identity zoom, use \(f_{-j^2e_1,j}\) with \(j\ge3\); its full image ball lies in the negative half-space and misses all supports. The two normalized embeddings are exactly \(Q^*L_\pi\) and \(L_\pi\). Uniform derivative and inverse-derivative bounds follow from \(|t'(r)|r\le2\) in every rescaled annulus, so the patched map remains quasiconformal. This repairs the intended nonuniqueness example. The main proof does not depend on the stated numerical choices, hence Minor.\par
\section*{Positive checks, boundary cases, and rejected concerns}
Sobolev algebra and exact forms. For degrees \(k,\ell\) with \(k+\ell\le n\), Hölder exponents obey \(k/n+\ell/n=(k+\ell)/n\). The product derivative has the same integrability because each exterior derivative is controlled in the corresponding degree’s \(L^{n/k}\) space. For \(1\le k<n\), \(n/k>1\), so weak compactness is available; the top degree is handled separately as a measure. The independent script checks 680 such exponent cases for \(2\le n\le16\), supplementing this general calculation.\par
For an exact \(k\)-form \(d\eta\), integration by parts against a compact test form and the bound on \(f^*\eta\) leave the factor \(A(f)^{-k/n} A(f)^{(k-1)/n}=A(f)^{-1/n}\). This tends to zero. Degree zero uses a constant harmonic representative, not an unavailable reflexive \(L^\infty\) argument. Compactness of the regularized Poincaré operator produces strongly convergent primitives on a finite harmonic basis. Combining one strong primitive with one weak closed pullback and integrating by parts proves the distributional wedge limit, including total degree \(n\). The general operator theorem remains an external analytic input, not an independently reproved monograph.\par
Top measure and injectivity. For a non-rational-homology-sphere target, a single nonzero dual pair yields an \(L^1\) density for the vague Jacobian limit. Nonnegativity follows from positive approximating measures. Absolute continuity makes the unit sphere a null set, so restriction to \(\overline B^n\) preserves the normalized mass one. Poincaré duality then makes a nonzero class pair to a nonzero multiple of that top class, proving injectivity. In the rational-homology-sphere case the intermediate cohomology vanishes and a normalized top form can be chosen directly; the argument does not assert the same Jacobian-limit regularity there.\par
For pointwise evaluation in Section 7, choose representatives of a finite homogeneous basis and extend linearly. Intersect the finitely many full-measure sets on which the basis multiplication identities hold, and choose a point where the top density is nonzero. There is no need to intersect uncountably many full-measure sets. Evaluation is then a graded homomorphism, injective by the same duality argument.\par
Growth and compactness dependencies. Bonk–Heinonen [Retained latex: \textbackslash{}cite[Theorem 1.11]\{BH\}] supplies a positive growth exponent depending on dimension and distortion for the required nonzero intermediate cohomology. Its exact primary statement was checked; the apparent concern about an unmentioned target-dependent exponent was rejected. The hunting step extracts a fixed doubling constant and unbounded normalized area.\par
The bounded-area argument in Section 9 uses a local bounded-energy bubbling statement, whereas the prominently cited Pankka–Souto theorem is phrased for closed sources and fixed degree. The local proof ingredients in [Retained latex: \textbackslash{}cite[Propositions 3.1--3.2]\{PS\}] were checked: uniform energy bounds the number of bad balls; an exhaustion and diagonal subsequence give finitely many interior concentration points; finite-energy removable singularities yield spherical bubbles. A nonconstant sphere-to-\(N\) map would have positive degree, forcing an injection of intermediate cohomology into that of the sphere. Thus a target which is not a rational homology sphere has no such bubbles. This supplies the local adaptation used here, subject to those external compactness/removability inputs.\par
Crucially, Corollary 1.4 and Theorem 9.1 really contain the closure \(\overline{\mathcal F_{K,D}(N)}\). The overbar is visible in the source PDF and author TeX, although lost by plain-text extraction. Shrinking maps converging to a constant therefore do not refute the printed statements. This candidate concern was explicitly rejected.\par
Examples and four-dimensional algebra. For the sphere construction in Section 8, the tripling map fixes the equator as required for gluing. Each small inserted bubble has mass \(\operatorname{vol}(S^n)+O(j^{-n})\), the normalized weights are \(1/j+O(j^{-n-1})\), and the mass outside the small balls tends to zero. Uniform continuity and the midpoint Riemann sum give the claimed half-length measure on the segment. This is consistent with the exclusion of rational homology spheres in the absolute-continuity theorem.\par
On \(\bigwedge^{2m}\mathbb R^{4m}\), interpret the top-degree pairing as the coefficient of the fixed oriented volume element. Each complementary pair of basis monomials gives a \(2\)-by-\(2\) block with eigenvalues \(+1,-1\). There are \(\frac12\binom{4m}{2m}\) such blocks, proving the claimed signature bound. Exact checks give \((3,3)\), \((35,35)\), \((462,462)\), and \((6435,6435)\) for dimensions \(4,8,12,16\). The integral over all of \(\mathbb R^{4m}\) printed in Appendix A should be read as this algebraic evaluation (or integration over a unit-volume cube); a literal improper integral of a nonzero constant top form would not be finite. This notational repair leaves the signature argument intact and is not counted as a separate obstruction.\par
The PL argument in Appendix B has a locally evident composition-order typo, \(p\circ\iota\) in place of \(\iota\circ p\). With the corrected order, finite-simplex derivative bounds, cancellation of opposite boundary faces in Stokes’ formula, and a positive minimum Jacobian on each compact simplex give the asserted Sobolev regularity and finite distortion. This verification preserves the compatibility hypotheses emphasized in F02.\par
\section*{Conclusion, repairs, and adoption}
The findings distinguish a disproved proof step, an unproved extension of scope, and a repaired auxiliary example. No claim of complete correctness follows from the positive checks.\par
Item | Judgment | Repair or positive evidence | Adoption / remaining work\par
F01 | Major | Exact decomposition without division | Local repair proved; core source step invalid\par
F02 | Major | Standard representatives; necessary bound | Narrow quantifier or supply transfer theorem\par
F03 | Minor | Start at \(j=3\); use quarter-radius zoom | Corrected nonuniqueness example verified\par
Positive | Adopt locally | Weak wedge, duality, closure, signature | External analytic/topological inputs remain\par
Version | Limited | v2 and later-version metadata checked | Later accepted and journal text unverified\par
\section*{Supporting checks}
Input PDF SHA256:\par
> 00000TOKEN\par
Manifest SHA256:\par
> 00001TOKEN\par
Author TeX SHA256:\par
> 00002TOKEN\par
> 00004TOKEN\par
It returned exit code zero and recorded 21 exact rational metric checks, explicit exterior-form identities, eight normalization checks, supplementary periodic-mean checks, 20 twist-Jacobian checks, overlapping-ball tests and repaired-domain checks, middle-form signature enumeration, and 680 Hölder-exponent checks. The general proofs above do not depend on a finite numerical grid being exhaustive. No author code was run. The accepted-manuscript retrieval failure is retained in `dependencies/HP\_accepted\_Nov2024\_receipt.json`; it is not concealed as a successful version comparison.\par
The exact report is `reports/0009.tex`, compiled to `reports/Pdf/0009.pdf`. The source/PDF checksums are external to the PDF to avoid self-reference. The page-render and actual visual-check records bind every checked report page to that output.\par
Cache and temporary use were task-local and bounded, with both disks checked above the required free-space floor.\par
\section*{Supporting checks}
The visible reviewing assistant performed all scientific reading, proofs, independent calculations, static checks, compilation, and actual report-page inspection. No scientific work was delegated. General quasiregular mapping theory, Hodge theory, the compact homotopy-operator theorem, bubbling/removability results, and the full Freedman–Donaldson classification were not reproved from foundations. Their interfaces were examined where used. The report is a systematic initial review with explicit evidence and limits, not a correctness certificate.\par
The dispatcher’s actual observed interface was Minhao Liu’s Group — Academic Researchers, Chat mode, webpage 6Pro. This is attributed to the dispatcher. No model switch was requested or performed.\par
> 00016TOKEN\par
9 S. Heikkilä and P. Pankka, De Rham algebras of closed quasiregularly elliptic manifolds are Euclidean, arXiv:2302.11440v2 (2023), exact reviewed version. https://arxiv.org/abs/2302.11440v2 Annals of Mathematics, DOI 10.4007/annals.2025.201.2.3, publisher identity/revision metadata only. https://annals.math.princeton.edu/2025/201-2/p03 R. Piergallini and D. Zuddas, Branched coverings of \(\mathbb{CP}^2\) and other basic 4-manifolds, arXiv:1707.03667v4 (2021), Theorem 1.3, Remark 1.4, Corollary 1.5. https://arxiv.org/abs/1707.03667v4 P. Pankka and J. Souto, Bubbling of quasiregular maps, arXiv:1904.00885v1 (2019), especially local Propositions 3.1–3.2. https://arxiv.org/abs/1904.00885v1 M. Bonk and J. Heinonen, Quasiregular mappings and cohomology, Acta Math. 186 (2001), 219–238, Theorem 1.11, primary PDF inspected. https://doi.org/10.1007/BF02401840\par
\end{document}
